% EXTRACTED from v1.1 manuscript (tag v1.1-verified), Part VIII,
% "The simply typed lambda calculus", lines 12033--12071.  Type
% metavariables are Greek in the handbook so they do not collide with the
% Part-7 formula metavariables used by the proof-theory chapters.
\calculussheet{sheet:stlc}{Simply Typed Lambda Calculus $\mathsf{STLC}$}{Types form the least set containing every $\mathsf b\in\mathsf{Base}$ and closed under $(\alpha,\beta)\mapsto\alpha\mathbin{\to_{\mathrm{ty}}}\beta$; terms form the least set containing variables and closed under typed abstraction $\lambda x{:}\alpha.M$ and application $MN$; the locally declared typing judgment is $\Theta_{\mathrm{ty}}\derives_{\mathsf{STLC}}M{:}\alpha$, where $\Theta_{\mathrm{ty}}$ is a finite list assigning distinct term variables to types; $\to_{\mathrm{ty}}$ is the registered metalinguistic type former, not object-language implication; term substitution is capture avoiding; this arrow-only calculus has no products, sums, polymorphism, or dependent types.}
\hbrule{rule:stlc-var}{Var}{x{:}\alpha\in\Theta_{\mathrm{ty}}}{\Theta_{\mathrm{ty}}\derives_{\mathsf{STLC}}x{:}\alpha}{\none}
\hbrule{rule:stlc-abs}{$\to_{\mathrm{ty}} I$}{\Theta_{\mathrm{ty}},x{:}\alpha\derives_{\mathsf{STLC}}M{:}\beta}{\Theta_{\mathrm{ty}}\derives_{\mathsf{STLC}}\lambda x{:}\alpha.M{:}\alpha\mathbin{\to_{\mathrm{ty}}}\beta}{$x$ is not declared in $\Theta_{\mathrm{ty}}$; the bound variable may be renamed consistently, and typing is stable under that alpha-renaming}
\hbrule{rule:stlc-app}{$\to_{\mathrm{ty}} E$}{\Theta_{\mathrm{ty}}\derives_{\mathsf{STLC}}M{:}\alpha\mathbin{\to_{\mathrm{ty}}}\beta\qquad\Theta_{\mathrm{ty}}\derives_{\mathsf{STLC}}N{:}\alpha}{\Theta_{\mathrm{ty}}\derives_{\mathsf{STLC}}MN{:}\beta}{\none}
\hbrule{rule:stlc-beta}{$\beta$}{}{(\lambda x{:}\alpha.M)N\longrightarrow_{\beta}M[N/x]}{term substitution $M[N/x]$ is capture avoiding; compatible beta reduction is the closure of this contraction under term contexts}

For the corresponding arrow-only natural-deduction rules, read each base
type as an atom and the type arrow as implication.  The following is the
localized intuitionistic presentation of
\cite[\S1.3, pp.~15--17]{Pfenning1998LocalizedND}: $\Gamma$ is a finite list
of distinctly labelled reusable formula assumptions, not a list of typed
terms.  An assumption permission may be unused.  It is therefore not the
classical ND-CPL sheet's nonempty-discharge presentation.
\hbrule{rule:p14-stlc-nd-assumption}{Hyp}{}{\Gamma_1,u:A,\Gamma_2\vdash A}{the assumption label $u$ is declared once in the context}
\hbrule{rule:p14-stlc-nd-intro}{$\supset I^u$}{\Gamma,u:A\vdash B}{\Gamma\vdash A\supset B}{$u$ is fresh for $\Gamma$; discharge its declaration, whether used or unused}
\hbrule{rule:p14-stlc-nd-elim}{$\supset E$}{\Gamma\vdash A\supset B\qquad\Gamma\vdash A}{\Gamma\vdash B}{the formula context is shared, not split linearly}
Term variables decorate the Hyp occurrences, abstraction decorates
$\supset I$, and application decorates $\supset E$.  This displayed
correspondence does not erase the capture-avoidance or typing conditions
of the STLC rows, or supply a classical RAA term constructor.
